\documentclass[11pt]{article}
\usepackage[margin=1in]{geometry}
\usepackage{amsmath,amssymb,amsthm}
\usepackage{graphicx}
\usepackage[hidelinks]{hyperref}
\setlength{\emergencystretch}{2em}

\newtheorem{theorem}{Theorem}
\newtheorem{corollary}{Corollary}
\newtheorem{remark}{Remark}

\newcommand{\C}{\mathbb C}
\newcommand{\dist}{\operatorname{dist}}

\title{BCP Blockers Have Essential Hereditary Support}
\author{Run \texttt{fa\_banach\_001}, lane 19}
\date{June 17, 2026}

\begin{document}
\maketitle

\section*{Source Problem and Context}

The source paper is Jinghao Huang, Karimbergen Kudaybergenov, and Rui Liu,
\emph{The ball-covering property of von Neumann algebras and noncommutative
symmetric spaces}, arXiv:2311.14257.  It asks for a noncommutative
\(C^*\)-algebra version of the commutative \(C(K)\) BCP/UBCP criterion:

\begin{center}
\includegraphics[width=0.92\textwidth]{figures/open_problem_crop.png}
\end{center}

The lane packet
\[
\texttt{2311.14257\_bcp\_countable\_capture\_exact\_equivalence}
\]
proved that a nonzero normed space has BCP if and only if there is a
countable set \(C\) such that every unit vector is not an exact capture
direction for \(C\).  Equivalently, for every \(u\in S_X\) some \(c\in C\)
satisfies
\[
        \dist(c,\mathbb F u)<\|c\|.
\]
We call such a set a BCP blocker.

The earlier packet
\[
\texttt{2311.14257\_bcp\_forces\_faithful\_state\_essential\_blocker\_ideal}
\]
proved that the closed two-sided ideal generated by a BCP blocker is
essential.  The present packet strengthens this from the ideal level to the
hereditary-support level.

\section*{Statement}

\begin{theorem}[essential hereditary support of a BCP blocker]
\label{thm:main}
Let \(A\) be a nonzero \(C^*\)-algebra with BCP.  Let
\[
        C=\{c_1,c_2,\ldots\}\subset A
\]
be a countable BCP blocker, so for every \(u\in S_A\) there is \(n\) with
\[
        \dist(c_n,\C u)<\|c_n\|.
\]
Define
\[
        h
        =
        \sum_{n=1}^\infty
        2^{-n-2}\,
        \frac{c_n^*c_n+c_nc_n^*}{1+\|c_n\|^2}
        \in A_+ .
\]
Then the hereditary subalgebra
\[
        B=\overline{hAh}
\]
is essential in \(A\); equivalently its annihilator
\[
        \operatorname{Ann}(B)=\{a\in A:aB=Ba=\{0\}\}
\]
is zero.
\end{theorem}

\begin{corollary}
\label{cor:sigma}
Every \(C^*\)-algebra with BCP contains an essential \(\sigma\)-unital
hereditary subalgebra generated canonically by a countable BCP blocker.
\end{corollary}

\section*{Idea of the Proof}

The two-sided ideal generated by a countable set can be much larger than any
separable or \(\sigma\)-unital local object.  For example, a rank-one
projection in \(\mathcal B(H)\), with \(H\) nonseparable, generates
\(\mathcal K(H)\) as an ideal.

The hereditary subalgebra generated by the blocker squares is more rigid.
If it had a nonzero annihilator, choose a positive unit vector \(u\) in that
annihilator.  Since \(h\) dominates positive multiples of both
\(c_n^*c_n\) and \(c_nc_n^*\), the relations \(uh=hu=0\) force
\[
        c_nu=uc_n=0 \qquad(n\ge1).
\]
Then \(c_n\) and \(u\) live in orthogonal corners, so
\[
        \|c_n-\lambda u\|=\max\{\|c_n\|,|\lambda|\}
        \qquad(\lambda\in\C).
\]
Thus \(u\) captures all of \(C\), contradicting the blocker property.

\section*{Proof}

\begin{proof}[Proof of Theorem~\ref{thm:main}]
The series defining \(h\) converges in norm and gives a positive element of
\(A\).  For every \(n\), there is a constant \(\alpha_n>0\) such that
\[
        h\ge \alpha_n c_n^*c_n,
        \qquad
        h\ge \alpha_n c_nc_n^* .
\]

Suppose for contradiction that \(B=\overline{hAh}\) is not essential.  Then
its annihilator contains a nonzero positive element.  Choose
\[
        0\le u\in \operatorname{Ann}(B),
        \qquad
        \|u\|=1.
\]
Since \(h\in B\), we have
\[
        uh=hu=0.
\]
For a fixed \(n\), the domination above gives
\[
        0\le u c_n^*c_n u\le \alpha_n^{-1} uhu=0,
\]
and hence \(c_nu=0\).  Similarly,
\[
        0\le u c_nc_n^* u\le \alpha_n^{-1} uhu=0,
\]
so \(c_n^*u=0\), equivalently \(uc_n=0\).

Therefore \(c_nu=uc_n=0\) for every \(n\).  It follows that for every
\(\lambda\in\C\),
\[
\begin{aligned}
        (c_n-\lambda u)^*(c_n-\lambda u)
        &=
        c_n^*c_n+|\lambda|^2u^2,
\end{aligned}
\]
and the two positive summands are orthogonal.  Hence
\[
        \|c_n-\lambda u\|^2
        =
        \max\{\|c_n^*c_n\|,|\lambda|^2\|u^2\|\}
        =
        \max\{\|c_n\|^2,|\lambda|^2\}.
\]
Thus
\[
        \dist(c_n,\C u)=\|c_n\|
        \qquad(n\ge1).
\]
This says that the unit vector \(u\) captures the whole blocker \(C\), a
contradiction.  Therefore \(\operatorname{Ann}(B)=0\), so \(B\) is essential.
\end{proof}

\begin{proof}[Proof of Corollary~\ref{cor:sigma}]
The hereditary subalgebra generated by a single positive element is
\(\sigma\)-unital; for instance, it has a countable approximate identity
obtained from continuous functional calculus of \(h\).  Apply
Theorem~\ref{thm:main}.
\end{proof}

\section*{Relation to the Remaining BCP/UBCP Wall}

This theorem does not prove UBCP.  Essential \(\sigma\)-unital hereditary
support is still weaker than a countable local coefficient model.  The
separable-essential-ideal packet proves UBCP when the essential local object
is separable, but the present \(B=\overline{hAh}\) can be nonseparable.

The theorem nevertheless sharpens the structural route.  A hypothetical
\(C^*\)-algebra with BCP but not UBCP cannot merely have a diffuse or remote
BCP blocker: the blocker produces a concrete positive element whose
hereditary support meets every nonzero ideal.  The remaining question is:

\begin{center}
\emph{Does an essential hereditary support coming from a BCP blocker contain
enough countable local coefficient structure to build a uniform ball cover?}
\end{center}

If yes, BCP implies UBCP for \(C^*\)-algebras.  If no, the failure should
produce one of the known countable-capture mechanisms.

\section*{Novelty and Search Bounds}

This packet was produced during the second consolidation of the lane-19
arXiv:2311.14257 work.  The local run indexes were searched for essential
blocker ideals, separable right-support nets, hereditary support, essential
hereditary subalgebras, and BCP/countable capture.  The previous
essential-blocker packet recorded essentiality of the ideal generated by the
blocker, but not this hereditary-support strengthening.

\section*{Human Review Focus}

Reviewers should check:
\begin{enumerate}
    \item that \(h\) dominates positive multiples of both \(c_n^*c_n\) and
          \(c_nc_n^*\);
    \item that a nonzero annihilator element \(u\) of \(\overline{hAh}\)
          satisfies \(uh=hu=0\);
    \item the implications \(hu=uh=0\Rightarrow c_nu=uc_n=0\);
    \item the orthogonal positive-summand norm calculation;
    \item the scope limitation: essential \(\sigma\)-unital hereditary support
          is not yet the same as UBCP.
\end{enumerate}

\begin{thebibliography}{9}

\bibitem{HKL}
Jinghao Huang, Karimbergen Kudaybergenov, and Rui Liu,
\emph{The ball-covering property of von Neumann algebras and noncommutative
symmetric spaces}, arXiv:2311.14257, 2023.

\bibitem{Blackadar}
Bruce Blackadar,
\emph{Operator Algebras: Theory of \(C^*\)-Algebras and Von Neumann
Algebras},
Encyclopaedia of Mathematical Sciences 122, Springer, 2006.

\bibitem{lane19capture}
Run \texttt{fa\_banach\_001}, lane 19,
\emph{BCP and Countable Capture Are Exact Opposites},
solution packet
\texttt{2311.14257\_bcp\_countable\_capture\_exact\_equivalence}, 2026.

\bibitem{lane19essential}
Run \texttt{fa\_banach\_001}, lane 19,
\emph{BCP Forces a Faithful State and an Essential Blocker Ideal},
solution packet
\texttt{2311.14257\_bcp\_forces\_faithful\_state\_essential\_blocker\_ideal},
2026.

\end{thebibliography}

\end{document}
